\section{Two fixed forward fields determine the table and launch law}
\label{sec:two-field-rigidity}
Fix a unit vector $e$, a length $t>0$, and write $a=te$.
The only commands in this section are $a$ and $-a$, both made with
the same unknown stationary launch law.  In particular, no angular
or length derivative belongs to the datum.  The two observed means
are the bounded measurable functions
\begin{equation}\label{eq:two-field-data}
 F_+(x)=\int_A B_a(x+z)\dd\mu(z),
 \qquad
 F_-(x)=\int_A B_{-a}(x+z)\dd\mu(z),
 \qquad x\in\R^2.
\end{equation}
These are active spatial fields at one fixed positive command
length.  They use the original first-collision bit, including its
zero value at a solid start, and all attempted preparations remain
in the denominator.

Let $\mathcal O=\bigcup_{C\in\mathcal C}C$ be a nonempty, locally
finite configuration of compact convex bodies with nonempty
interiors, diameters at most $D$, and mutual distances at least
$d>0$.  Each body is strictly convex; no boundary smoothness or
positive curvature bound is imposed in this section.
The unknown stationary law $\mu$ is a compactly supported Borel
probability measure whose support $A$ is nonempty and convex, with
$\diam A\le\Delta$.  The support may be lower dimensional and the
law may have atoms or a singular part.
The only size comparison is
\begin{equation}\label{eq:two-field-gap}
                         t+\Delta<d.
\end{equation}
There is no lower obstacle-width assumption, and the position of
$A$ is unrestricted.  No density or smoothness assumption is made
on the law.  Equality of the two data fields may be understood
almost everywhere with respect to planar Lebesgue measure.
As before, $h_K$ denotes the support function of a
compact set $K$, equivalently of its convex hull, and $s(K)$ denotes
the Steiner point when $K$ is nonempty and convex, including the
lower-dimensional cases.  Supports of observed fields below mean
supports of the nonnegative measures $F_\pm(x)\dd x$.

\begin{lemma}[A finite prefix inverse for occupation]
\label{lem:two-field-prefix}
Put
\begin{equation}\label{eq:two-field-increment}
 r(x)=F_+(x)-F_-(x+a),\qquad
 M=\left\lfloor\frac{D+\Delta}{t}\right\rfloor+1.
\end{equation}
The occupation field is given by the finite formula
\begin{equation}\label{eq:two-field-prefix-inverse}
 v(x):=\int_A\one_{\mathcal O}(x+z)\dd\mu(z)
 =\max_{0\le m\le M}
       \left\{-\sum_{k=0}^{m-1}r(x+ka)\right\}.
\end{equation}
The empty sum is zero.  If arbitrary supplied values of $r$ on
these $M$ sites differ by at most $\epsilon$ from the exact values,
the value returned by the same formula differs from $v(x)$ by at
most $M\epsilon$.  In particular, perturbing each of the two mean
fields by at most $\epsilon$ changes the returned value by at most
$2M\epsilon$.  The formula holds pointwise for the physical mean
fields.  If those fields are supplied only as almost-everywhere
equivalence classes, it determines the corresponding class of $v$.
\end{lemma}
\begin{proof}
For every actual start $y$, convexity and $t<d$ give the
endpoint identity
\begin{equation}\label{eq:two-field-bit-identity}
 B_a(y)-B_{-a}(y+a)
                  =\one_{\mathcal O}(y+a)-\one_{\mathcal O}(y).
\end{equation}
Indeed, if precisely one endpoint is solid, the bit pointing toward
it is one and the other is zero.  If both endpoints are free, the
two bits agree.  If both are solid, they lie in the same convex
body and both bits vanish.  The closed-set convention includes
tangencies and boundary endpoints: a boundary solid start is
zero, while a swept boundary reached from a free start counts as
a collision.  Thus the identity remains true on those sets,
which may have positive probability under a singular law.
Integrating against the same stationary probability measure proves
\begin{equation}\label{eq:two-field-occupation-increment}
                         r(x)=v(x+a)-v(x).
\end{equation}
The nominal shift in $F_-(x+a)$ realizes the endpoint comparison
using only the fixed forward command $-a$ with the original law.

The field $v$ is bounded and measurable.  Positivity of the measure
on every neighborhood of every point in $A$ gives
\begin{equation}\label{eq:two-field-occupation-components}
 \operatorname{supp}(v(x)\dd x)=\bigcup_{C\in\mathcal C}P_C,
 \qquad P_C=C+(-A).
\end{equation}
Indeed, each individual occupation measure is the convolution
of $\one_C(x)\dd x$ with the reflected probability measure,
and the support of a convolution of two nonnegative compactly
supported measures is the sum of their supports.  The sets $P_C$
are also the closures of the components of $\{v>0\}$:
every point of $\operatorname{int}(C+(-A))=\operatorname{int}C+(-A)$
has positive occupation, and outside $C+(-A)$ the contribution
of $C$ is zero.  Boundary occupation values need not vanish.
These compact convex component closures have diameters at most
$D+\Delta$ and mutual gaps at least $d-\Delta>t$.
Every sum in \eqref{eq:two-field-prefix-inverse} telescopes to
$v(x)-v(x+ma)$ and is therefore at most $v(x)$.
If $v(x)=0$, the empty sum attains equality.  Otherwise $x$ is in
one component $P_C$.  Since $Mt>D+\Delta$, the ray
$x,x+a,\ldots,x+Ma$ exits that component by time $M$.
At its first exit it cannot land in another component, because
the component gap exceeds $t$.  Its occupation at that point is
zero, so the corresponding prefix attains equality.
Finally, each prefix sum changes by at most $m\epsilon$, and
taking a maximum is nonexpansive in the maximum norm.  This proves
the stability bounds.
\end{proof}

The inverse is a finite directed version of the prefix construction
in the earlier local experiment.  It is useful here because the
two commanded directions are retained as separate channels.
It has a fixed finite spatial range $Mt$ and does not require a
limiting walk, an angularly averaged forcing, or prior knowledge of
the unknown positive components.

The bounded measurability of the fields is sufficient here: each
prefix uses only finitely many fixed translations, so changing
the supplied fields on null sets changes the recovered occupation
only on a null set.

\subsection{Measure supports and singular launch laws}
\label{subsec:measure-support}
For a nonnegative locally integrable function $f$, the convention
used here is
\[
 \operatorname{supp}(f(x)\dd x)
 =\left\{x:\int_{B(x,r)}f(y)\dd y>0
                         \text{ for every }r>0\right\}.
\]
This set depends only on the almost-everywhere class of $f$.
It need not be the pointwise positivity set.  We write
$\check\mu=(-\operatorname{id})_\#\mu$.

\begin{lemma}[Positive convolution supports]
\label{lem:positive-convolution-support}
Let $\alpha,\beta$ be nonzero, nonnegative, compactly supported
finite Borel measures on $\R^2$.  Then
\begin{equation}\label{eq:positive-convolution-support}
 \operatorname{supp}(\alpha*\beta)
            =\operatorname{supp}\alpha+\operatorname{supp}\beta.
\end{equation}
If bounded measurable sets $E,E'$ differ by a planar null set,
then $\one_E*\beta=\one_{E'}*\beta$ almost everywhere,
including when $\beta$ is singular.
\end{lemma}
\begin{proof}
Put $U=\operatorname{supp}\alpha$ and
$V=\operatorname{supp}\beta$.  The convolution is carried by the
compact set $U+V$.  Conversely, if $x=p+q$ with $p\in U$ and
$q\in V$, then, for every $r>0$,
\[
 (\alpha*\beta)(B(x,r))
 \ge \alpha(B(p,r/3))\beta(B(q,r/3))>0.
\]
This proves \eqref{eq:positive-convolution-support}.
Tonelli's theorem identifies the convolution
$(\one_E(x)\dd x)*\beta$ with the measure having density
$\one_E*\beta$.  It also gives
\[
 \int_{\R^2}|\one_E*\beta-\one_{E'}*\beta|\dd x
 \le \beta(\R^2)\,\area(E\mathbin\triangle E')=0.
\]
Neither argument requires a density for either measure.
\end{proof}

\begin{proposition}[Observable components under an arbitrary compact law]
\label{prop:singular-support-components}
Under \eqref{eq:two-field-gap}, let
$v_C(x)=\int\one_C(x+z)\dd\mu(z)$, and let
$f_{\pm,C}(x)=\int B^{C}_{\pm a}(x+z)\dd\mu(z)$,
where $B^{C}$ is the collision bit for the single obstacle $C$.
Use coordinates $y=ue+we^\perp$ and write the sections of $C$ as
$[\ell_C(w),u_C(w)]$, with $w\in I_C=[w_-,w_+]$.
Define the closed arcs and strips by
\begin{align*}
 \Gamma_{+,C}&=\{\ell_C(w)e+we^\perp:w\in I_C\},&
 E_{+,C}&=\Gamma_{+,C}+[-t,0]e,\\
 \Gamma_{-,C}&=\{u_C(w)e+we^\perp:w\in I_C\},&
 E_{-,C}&=\Gamma_{-,C}+[0,t]e.
\end{align*}
Then
\begin{equation}\label{eq:singular-individual-supports}
 \operatorname{supp}(v_C(x)\dd x)=P_C:=C+(-A),\qquad
 \operatorname{supp}(f_{\pm,C}(x)\dd x)
                         =K_{\pm,C}:=E_{\pm,C}+(-A).
\end{equation}
The connected components of the occupation support are precisely
the $P_C$; those of the two collision supports are, respectively,
the $K_{+,C}$ and the $K_{-,C}$.  Moreover,
\begin{align}
 \dist(P_C,P_{C'})&\ge d-\Delta,\label{eq:singular-occupation-gap}\\
 \dist(K_{\pm,C},K_{\pm,C'})&\ge d-t-\Delta,
       \qquad C\ne C',\label{eq:singular-collision-gap}\\
 K_{\pm,C}\cap P_C&\ne\varnothing,\qquad
 \dist(K_{\pm,C},P_{C'})\ge d-t-\Delta
       \quad(C\ne C').\label{eq:singular-support-matching}
\end{align}
Thus intersection with the occupation support determines the
correspondence between the two collision component families.
These statements apply to point, segment, atomic and singular
continuous launch laws with the stated convex support.
\end{proposition}
\begin{proof}
The section functions $\ell_C$ and $u_C$ are convex and concave,
respectively, on $I_C$.  They are continuous in its interior.
Strict convexity makes each extreme section a singleton.  Any
convergent subsequence of section endpoints as $w\to w_\pm$
lies in that singleton, by compactness of $C$.  Both section
functions are consequently continuous on $I_C$.

For the positive command, the single-obstacle collision-start set
is exactly
\[
 D_{+,C}=\{(\ell_C(w)-s)e+we^\perp:
                         w\in I_C,\ 0<s\le t\}.
\]
The endpoint $s=0$ is a solid start and is excluded; $s=t$ is
a collision at the terminal endpoint and is included.  At
$w=w_\pm$ the same formula includes a tangential encounter.
The open set
\[
 G_{+,C}=\{(\ell_C(w)-s)e+we^\perp:
                    w_-<w<w_+,\ 0<s<t\}
\]
has closure $E_{+,C}$.  Indeed, the parametrization is a
homeomorphism from $I_C\times[0,t]$ to $E_{+,C}$, and its
restriction to the open rectangle has open image.  The sets
$D_{+,C}$ and $G_{+,C}$ differ only on two graphs and two
transverse endpoint segments, all of which are planar null sets
by one-dimensional slicing.  The same description, with
$u_C(w)+s$ in place of $\ell_C(w)-s$, applies to the negative
command.  In particular,
\[
 \operatorname{supp}(\one_{G_{\pm,C}}(x)\dd x)=E_{\pm,C}.
\]
Every neighborhood of a point of $E_{\pm,C}$ meets its open
interior in positive area, including at a tangent endpoint.

Since $C$ is the closure of its interior,
$\operatorname{supp}(\one_C(x)\dd x)=C$.  Apply
Lemma~\ref{lem:positive-convolution-support} to these three
Lebesgue measures and to $\check\mu$, whose support is $-A$.
The null-set assertion in that lemma passes from $G_{\pm,C}$
to the physical collision bits after convolution and proves
\eqref{eq:singular-individual-supports}.  This passage is an
almost-everywhere assertion; it does not discard the prescribed
pointwise boundary convention for the physical means.

For completeness, the occupation is positive at every point of
$\operatorname{int}P_C$.  The convex-set identity needed here is
\begin{equation}\label{eq:singular-interior-sum}
             \operatorname{int}(C+(-A))
                         =\operatorname{int}C+(-A).
\end{equation}
To prove the nontrivial inclusion, fix $c_0\in\operatorname{int}C$
and $a_0\in A$.  If $x\in\operatorname{int}(C-A)$, then for
small $\lambda>0$ the point
$(x-\lambda(c_0-a_0))/(1-\lambda)$ belongs to $C-A$.
Write it as $c-a$.  The identity
\[
 x=((1-\lambda)c+\lambda c_0)
                  -((1-\lambda)a+\lambda a_0)
\]
places $x$ in $\operatorname{int}C-A$; the reverse inclusion
follows because the latter set is a union of open translates.
If $x=c-a$ with $c\in\operatorname{int}C$ and $a\in A$, a
sufficiently small neighborhood of $a$ is mapped by $z\mapsto x+z$
into $C$.  That neighborhood has positive $\mu$-mass, so
$v_C(x)>0$.  Outside $P_C$ the occupation is zero.  Its positive
set therefore contains $\operatorname{int}P_C$ and is contained
in $P_C$; it is path connected, since a segment joining any of
its points to an interior point stays in the interior except
possibly at its first endpoint.  Its closure is $P_C$.

Each $E_{\pm,C}$ is path connected by its rectangle
parametrization.  Convexity of $A$ makes $K_{\pm,C}$ path
connected as well.  The estimates
\eqref{eq:singular-occupation-gap}--\eqref{eq:singular-support-matching}
follow by subtracting two representations of their points:
obstacle points have distance at least $d$, the difference of
two same-sign longitudinal shifts has length at most $t$, and
the difference of two footprint shifts has length at most
$\Delta$.  The nonempty intersection follows from
\[
 \Gamma_{\pm,C}+(-A)\subset K_{\pm,C}\cap P_C.
\]

All expanded families remain locally finite: a member meeting a
fixed compact set comes from an obstacle meeting its enlargement
by the fixed bounded set $A+[-t,t]e$.  Since $t<d$, a segment
cannot meet two different obstacles or start in one obstacle and
reach another.  Thus $v=\sum_Cv_C$ and
$F_\pm=\sum_C f_{\pm,C}$, with the physical boundary convention.
Their measure supports are the locally finite unions in
\eqref{eq:singular-individual-supports}.  The positive gaps and
connectedness just proved identify their connected components
and establish the matching assertion.
\end{proof}

For example, for any nonempty compact convex $A$, choose a sequence
$(a_j)_{j\ge1}$ dense in $A$ and put
$\mu=\sum_{j\ge1}2^{-j}\delta_{a_j}$.
Its topological support is $A$, and the proposition applies to this
purely atomic probability. This includes convex supports of dimension
zero, one or two and illustrates why the measure-support convention
is used throughout.

\begin{corollary}[Almost-everywhere data and the period group]
\label{cor:singular-support-ae-periods}
The almost-everywhere classes of $F_+,F_-$ determine the same
occupation and collision support components, with the same
matching, as their physical representatives.  If
\[
 \operatorname{Per}_{\mathrm{ae}}(F_+,F_-)
 =\{w:F_\pm(\cdot+w)=F_\pm\text{ almost everywhere}\},
\]
then
\begin{equation}\label{eq:singular-support-periods}
 \operatorname{Per}_{\mathrm{ae}}(F_+,F_-)
 =\operatorname{Per}_{\mathrm{point}}(F_+,F_-)
 =\operatorname{Per}(\mathcal O).
\end{equation}
\end{corollary}
\begin{proof}
The finite prefix formula in Lemma~\ref{lem:two-field-prefix}
uses only finitely many fixed translations.  A change of the
fields on null sets changes its output only on a finite union
of translated null sets.  It therefore recovers the
almost-everywhere occupation class and hence its measure support.
The two original measure supports and their incidence relation
are likewise unchanged.  Proposition~\ref{prop:singular-support-components}
identifies the components without any choice of representatives.

If $w$ is an almost-everywhere common field period, the same
prefix formula gives $v(\cdot+w)=v$ almost everywhere.
Translation by $w$ permutes the components of its measure
support, so $P_C+w=P_{C'}$ for some $C'$.  Taking support
functions and canceling the common $h_{-A}$ gives
$h_C(u)+w\cdot u=h_{C'}(u)$ for every unit $u$.
Hence $C+w=C'$ and $\mathcal O+w=\mathcal O$.
Conversely, a period of $\mathcal O$ preserves every physical
collision bit pointwise, and integration against the same law
preserves every physical mean pointwise.  This proves
\eqref{eq:singular-support-periods}.
\end{proof}


\begin{lemma}[The two collision supports]
\label{lem:two-field-supports}
For each body $C$, let $\Gamma_{+,C}$ be its incoming boundary arc
for direction $e$, including its two tangent endpoints, and let
$\Gamma_{-,C}$ be the opposite incoming arc for direction $-e$.
Define
\begin{equation}\label{eq:two-field-strips}
 E_{+,C}=\Gamma_{+,C}+[-t,0]e,\qquad
 E_{-,C}=\Gamma_{-,C}+[0,t]e,
 \qquad K_{\pm,C}=E_{\pm,C}+(-A).
\end{equation}
The connected components of $\operatorname{supp}F_\pm$ are exactly
the sets $K_{\pm,C}$.  Within either sign, different components
have distance at least $d-\Delta-t>0$.
The component $K_{\pm,C}$ meets $P_C$ and is disjoint from all
$P_{C'}$ with $C'\ne C$.  Thus the known occupation components
match the two families of collision supports without an unknown
label correspondence.  Their support functions satisfy
\begin{align}
 h_{K_{+,C}}(u)
   &=h_{\Gamma_{+,C}}(u)+h_{-A}(u)+t(-u\cdot e)_+,
       \label{eq:two-field-plus-support}\\
 h_{K_{-,C}}(u)
   &=h_{\Gamma_{-,C}}(u)+h_{-A}(u)+t(u\cdot e)_+.
       \label{eq:two-field-minus-support}
\end{align}
\end{lemma}
\begin{proof}
Use coordinates $y=ue+w e^\perp$.  The section of $C$ at transverse
coordinate $w$ is an interval $[\ell_C(w),u_C(w)]$ for
$w$ in its projection interval $[w_-,w_+]$.
The incoming graph $\ell_C$ is convex and continuous, and the
outgoing graph $u_C$ is concave and continuous.  The free starts
whose segment first hits $C$ in direction $e$ form, up to null
sets, the strip
\[
 \{we^\perp+(\ell_C(w)-s)e:
                         w_-<w<w_+,\ 0<s<t\}.
\]
Its closure is $E_{+,C}$; the reverse direction gives $E_{-,C}$.
No such strip meets another obstacle because $t<d$.
The planar density of $F_\pm(x)\dd x$ is consequently the locally
finite sum of convolutions of the strip indicators with the
reflected measure.  The strips and collision-start sets differ
only on planar null sets; convolution with any finite measure
preserves that equality almost everywhere, by Fubini's theorem.
Each strip is connected, has nonempty interior, and is the
closure of that interior.  The support of its convolution with
the nonnegative measure is exactly its Minkowski sum with $-A$:
there is no cancellation, and both factors give positive mass
to every neighborhood of every point in their supports.
Thus each $K_{\pm,C}$ is a connected compact support.

For two points from different strips with the same sign, the
two longitudinal shifts differ by at most $t$, while two
footprint shifts differ by at most $\Delta$.  Their distance is
therefore at least $d-t-\Delta$.  The same estimate separates
$K_{\pm,C}$ from $P_{C'}$ when $C'\ne C$.
Since $\Gamma_{\pm,C}\subset C$,
$\Gamma_{\pm,C}+(-A)\subset K_{\pm,C}\cap P_C$.
Local finiteness and the positive gap now prove the support
component and matching assertions.  Finally, support functions
add under Minkowski addition, giving the two displayed identities.
\end{proof}

\subsection{Support arcs and angular curvature without boundary regularity}
\label{subsec:nonsmooth-support-curvature}
The incoming arcs can be treated through sections, without a
parametrization by the boundary normal.  This also distinguishes
strict convexity from smoothness in the subsequent measure
calculation.

\begin{lemma}[A section form of the arc identity]
\label{lem:nonsmooth-arc-support}
Let $C$ be a strictly convex planar body with nonempty interior.
Let $p_-,p_+$ be its unique contacts with the support lines of
normals $-e^\perp,e^\perp$, and put $L=[p_-,p_+]$.
For its closed incoming arcs $\Gamma_+,\Gamma_-$ one has
\begin{align}
 h_{\Gamma_+}(u)&=
 \begin{cases}h_C(u),&u\cdot e\le0,\\
                h_L(u),&u\cdot e\ge0,
 \end{cases}
 \label{eq:nonsmooth-plus-hemispheres}\\
 h_{\Gamma_-}(u)&=
 \begin{cases}h_L(u),&u\cdot e\le0,\\
                h_C(u),&u\cdot e\ge0.
 \end{cases}
 \label{eq:nonsmooth-minus-hemispheres}
\end{align}
The two alternatives agree when $u\cdot e=0$.  In particular,
\begin{equation}\label{eq:nonsmooth-arc-identity}
                  h_{\Gamma_+}+h_{\Gamma_-}=h_C+h_L.
\end{equation}
\end{lemma}
\begin{proof}
Write $u=\alpha e+\beta e^\perp$ and use the interval sections
$[\ell(w),r(w)]$ for $w\in[w_-,w_+]$.  Convexity of $C$ makes
$\ell$ convex and $r$ concave.  They are continuous on the
closed projection interval: interior continuity follows from
convexity, and endpoint continuity follows from compactness and
the uniqueness of the extreme sections.  The endpoints of both
graphs are $p_-,p_+$.

If $\alpha\le0$, maximizing $\alpha s+\beta w$ over each section
chooses $s=\ell(w)$.  Consequently
$h_{\Gamma_+}(u)=h_C(u)$.  If $\alpha\ge0$, the function
$w\mapsto\alpha\ell(w)+\beta w$ is convex.  Its value at an
interior point is at most the larger endpoint value, by the
defining convexity inequality.  Its maximum is therefore
$\max\{u\cdot p_-,u\cdot p_+\}=h_L(u)$.
For the other arc, $\alpha\ge0$ selects the right endpoint of
each section and gives $h_C$; when $\alpha\le0$ the function
$\alpha r(w)+\beta w$ is convex and gives $h_L$.
At $\alpha=0$ both maxima depend only on the projection interval.
This proves all the formulas, including directions supporting
a corner and the two transverse directions.
\end{proof}

We identify the angular circle with $\R/(2\pi\Z)$, put
$n_\phi=(\cos\phi,\sin\phi)$ and
$\tau_\phi=n_\phi^\perp$, and use the planar support convention
\begin{equation}\label{eq:nonsmooth-area-convention}
 \int\psi\dd\mathsf S_K
     =\int_0^{2\pi}h_K(n_\phi)
                      (\psi''(\phi)+\psi(\phi))\dd\phi,
 \qquad \psi\in C^\infty(\R/(2\pi\Z)).
\end{equation}
Thus angular differentiation is distributional.  The following
proof supplies both positivity and the precise atom convention
for this surface-area measure.

\begin{lemma}[Surface-area atoms and exposed segments]
\label{lem:nonsmooth-surface-atoms}
For a compact convex planar body $K$ with nonempty interior,
\eqref{eq:nonsmooth-area-convention} defines a finite nonnegative
measure.  Let
\[
 F(K,n_\phi)=\{x\in K:x\cdot n_\phi=h_K(n_\phi)\}
\]
be its exposed face.  Writing $h(\phi)=h_K(n_\phi)$, its
one-sided derivatives satisfy
\begin{equation}\label{eq:nonsmooth-face-derivatives}
 h'_-(\phi)=\min_{x\in F(K,n_\phi)}x\cdot\tau_\phi,
 \qquad
 h'_+(\phi)=\max_{x\in F(K,n_\phi)}x\cdot\tau_\phi.
\end{equation}
Consequently
\begin{equation}\label{eq:nonsmooth-atom-face-length}
 \mathsf S_K(\{n_\phi\})
       =h'_+(\phi)-h'_-(\phi)
       =\operatorname{length}F(K,n_\phi).
\end{equation}
In particular, $\mathsf S_K$ is nonatomic if and only if $K$ is
strictly convex.  Boundary corners and a singular continuous
part of $\mathsf S_K$ are permitted in this equivalence.
\end{lemma}
\begin{proof}
Choose $R$ with $K\subset B(0,R)$.  The support function is
Lipschitz in $\phi$.  Each function
$x\cdot n_\phi+R\phi^2/2$, $x\in K$, is convex on the real
line, since its second derivative is $R-x\cdot n_\phi\ge0$.
Their supremum $h(\phi)+R\phi^2/2$ is convex.  Hence $h$ has
one-sided derivatives of locally bounded variation and its
distributional second derivative is a locally finite signed
measure.  Periodicity gives a finite signed measure on the circle.

If $p$ is a support point at $\phi$, then, for $|b|<\pi/2$,
\[
 h(\phi+b)+h(\phi-b)
 \ge p\cdot(n_{\phi+b}+n_{\phi-b})
 =2\cos b\,h(\phi).
\]
Multiply by a nonnegative smooth periodic $\psi$, integrate,
translate the two shifted integrals, and divide by $b^2$.
Letting $b\to0$ yields
$\int h(\psi''+\psi)\ge0$.  Thus the signed measure
$h''+h\dd\phi$ is nonnegative and is the measure in
\eqref{eq:nonsmooth-area-convention}.

For $b>0$ and $p\in F(K,n_\phi)$, the support inequality gives
\[
 \frac{h(\phi+b)-h(\phi)}{b}
       \ge p\cdot\frac{n_{\phi+b}-n_\phi}{b}.
\]
Conversely, choose $p_b\in F(K,n_{\phi+b})$.  The same
quotient is at most
$p_b\cdot(n_{\phi+b}-n_\phi)/b$.
Every subsequential limit of $p_b$ belongs to $F(K,n_\phi)$,
by compactness and continuity of $h$.  Taking lower and upper
limits proves the right-derivative formula in
\eqref{eq:nonsmooth-face-derivatives}; negative $b$ gives the
left-derivative formula.  The atom of the Stieltjes measure
$h''$ at $\phi$ is the jump $h'_+(\phi)-h'_-(\phi)$, whereas
$h\dd\phi$ has no atoms.  The face is a segment on the line
$x\cdot n_\phi=h(\phi)$, so its length is exactly the difference
between the two extremal $\tau_\phi$ coordinates.  This proves
\eqref{eq:nonsmooth-atom-face-length}.

The geometric interpretation of the measure follows in the same
coordinates.  The right-continuous support contact
$x_+(\phi)=h(\phi)n_\phi+h'_+(\phi)\tau_\phi$ selects the
positive tangent endpoint of each exposed face.  As the normal
turns once, these contacts traverse the convex boundary in order;
each jump is completed by the corresponding exposed segment.
The distributional product rule gives
\[
                      D x_+=\tau_\phi\dd\mathsf S_K.
\]
Since $|\tau_\phi|=1$ and $\mathsf S_K\ge0$, the variation
of this boundary parametrization is $\mathsf S_K$.  It assigns
boundary arclength to its outer normal, with a whole exposed
segment assigned to its fixed normal.  This is the usual planar
surface-area measure in the convention
\eqref{eq:nonsmooth-area-convention}.

A strictly convex body has only singleton exposed faces and
hence no atoms.  Conversely, any nontrivial boundary segment
lies in an exposed face: a supporting line at an interior point
of the segment contains both endpoints.  Such a segment would
give a positive atom.  At a corner of a strictly convex body,
the normal may vary over an interval while the support point
remains a fixed point $p$.  On the interior of that interval
$h(\phi)=p\cdot n_\phi$, so $h''+h=0$ there; its endpoints
also have zero atomic mass by
\eqref{eq:nonsmooth-atom-face-length}.  None of these arguments
excludes a singular continuous part of the measure.
\end{proof}

These are the planar support-measure conventions of
\cite[Chapter~4]{Schneider}.  Signed support differences and the
separate atomic, absolutely continuous and singular continuous
parts of their length measures are treated in
\cite[Section~4]{MartinezMaure}.  Formula
\eqref{eq:nonsmooth-area-convention} retains all three parts;
an almost-everywhere second derivative alone would not do so.

\begin{lemma}[Jordan recovery of the contact chord]
\label{lem:nonsmooth-jordan-chord}
For the body and chord in Lemma~\ref{lem:nonsmooth-arc-support},
put $H=h_C-h_L$, $m=(p_-+p_+)/2$, and
$\ell=|p_+-p_-|$.  If $n$ is either unit normal of the chord,
then the Jordan decomposition of $\sigma_H=(\partial_\phi^2+1)H$
is
\begin{equation}\label{eq:nonsmooth-jordan-parts}
 \sigma_H^+=\mathsf S_C,\qquad
 \sigma_H^-=\ell(\delta_n+\delta_{-n}).
\end{equation}
It determines the centered chord
$L^0=[-\ell n^\perp/2,\ell n^\perp/2]$, and
\begin{equation}\label{eq:nonsmooth-jordan-recovery}
                     H+h_{L^0}=h_{C-m}.
\end{equation}
Exactly one of the two chord normals has positive scalar product
with $e$.
\end{lemma}
\begin{proof}
The projection interval has positive length, so $\ell>0$.
Writing $q=(p_+-p_-)/\ell$, the chord support is
\[
 h_L(n_\phi)=m\cdot n_\phi+\frac\ell2|q\cdot n_\phi|.
\]
The translation term is annihilated by $\partial_\phi^2+1$.
The second term solves $h''+h=0$ away from the two normals
$n,-n$, and its first derivative jumps upward by $\ell$ at
each.  Hence
\[
 \sigma_H=\mathsf S_C-\ell(\delta_n+\delta_{-n}).
\]
Lemma~\ref{lem:nonsmooth-surface-atoms} shows that $\mathsf S_C$
assigns zero mass to these two points.  The two nonnegative
measures on the right are therefore mutually singular, so they
are exactly the positive and negative Jordan parts.  The
locations and equal masses of the negative atoms determine
$L^0$, independent of the choice between $n$ and $-n$.
Since $h_L=h_{L^0}+m\cdot u$, adding $h_{L^0}$ to $H$ gives
\eqref{eq:nonsmooth-jordan-recovery}.  Finally,
$(p_+-p_-)\cdot e^\perp=w_+-w_->0$; a chord normal therefore
has nonzero scalar product with $e$, and its sign selects exactly
one of the two normals.  The chord need not be perpendicular to
the command direction.
\end{proof}


\begin{lemma}[Calibration by a pair of supports]
\label{lem:two-field-calibration}
Let $P=P_C$ and let $K_+,K_-$ be its matched collision components.
The observable function
\begin{equation}\label{eq:two-field-calibration-function}
 H_C(u)=2h_P(u)-h_{K_+}(u)-h_{K_-}(u)+t|u\cdot e|
\end{equation}
determines $C$ up to translation.  More precisely, let $p_+,p_-$
be the two points of $C$ with outward normals $e^\perp,-e^\perp$,
and let $L_C=[p_-,p_+]$.  Then
\begin{equation}\label{eq:two-field-chord-cancellation}
                         H_C=h_C-h_{L_C}.
\end{equation}
Writing $u=n_\phi$, $\ell=|p_+-p_-|>0$, and choosing the unit
normal $n$ of this chord with $n\cdot e>0$, one has, as signed
measures on the angular circle,
\begin{equation}\label{eq:two-field-curvature-measure}
 (\partial_\phi^2+1)H_C
    =\mathsf S_C-\ell(\delta_n+\delta_{-n}),
\end{equation}
Here $\mathsf S_C=h_C''+h_C$ is the planar surface area measure,
which is nonnegative and has no atoms because $C$ is strictly
convex.  The negative part therefore determines the centered segment
$L_C^0=[-\ell n^\perp/2,\ell n^\perp/2]$.  The function
$H_C+h_{L_C^0}$ is the support function of $C-s(L_C)$.
\end{lemma}
\begin{proof}
On the hemisphere $u\cdot e\le0$, the support point of $C$ in
direction $u$ lies on $\Gamma_{+,C}$, so
$h_{\Gamma_{+,C}}(u)=h_C(u)$.  On the opposite hemisphere,
the function
\[
 w\longmapsto (u\cdot e)\ell_C(w)+(u\cdot e^\perp)w
\]
is convex.  Its maximum is attained at a projection endpoint,
so $h_{\Gamma_{+,C}}(u)=h_{L_C}(u)$ there.
For the reverse arc the two hemispheres are interchanged.
Thus, including their endpoints,
\begin{equation}\label{eq:two-field-arc-sum}
              h_{\Gamma_{+,C}}+h_{\Gamma_{-,C}}
                                      =h_C+h_{L_C}.
\end{equation}
Since $h_P=h_C+h_{-A}$, substitution of
\eqref{eq:two-field-plus-support} and
\eqref{eq:two-field-minus-support} into
\eqref{eq:two-field-calibration-function} proves
\eqref{eq:two-field-chord-cancellation}.  In particular, the unknown
footprint cancels exactly before differentiation is used.

For a planar convex body, $h_C''+h_C$ is the finite nonnegative
surface area measure.  Its atom at a normal direction has mass
equal to the length of the exposed boundary segment with that
normal.  In the strictly convex case it therefore has no atoms.
Equivalently, strict convexity gives a unique support point in
every direction; compactness makes that point continuous in the
direction, and differentiation of the support maximum gives
$h_C\in C^1$, so its first derivative has no jumps.  These are
the usual planar support-measure identities; see \cite{Schneider}. Signed differences of planar surface area
measures and their uniqueness up to translation are treated in
\cite[Section~4]{MartinezMaure}. Here the collision identity
\eqref{eq:two-field-chord-cancellation} supplies the particular
signed support whose atomic part can be separated.
A segment of length
$\ell$ has angular curvature measure
$\ell(\delta_n+\delta_{-n})$.  To check the coefficient directly,
after subtracting the midpoint translation its support function is
$\frac\ell2|n^\perp\cdot n_\phi|$: its first derivative has a
jump of $\ell$ at each of the two normal directions and it solves
$h''+h=0$ between them.
This proves \eqref{eq:two-field-curvature-measure}.
The nonnegative nonatomic measure and the two negative atoms are
mutually singular, so they are precisely the positive and negative
parts of this signed measure.  The two endpoints of the segment
have distinct transverse coordinates, hence $\ell>0$ and exactly
one of its normals has positive dot product with $e$.
Its length and unoriented direction are thus recovered from the
negative part.  Subtracting the segment midpoint from
\eqref{eq:two-field-chord-cancellation} gives the last assertion.
\end{proof}

The differentiation in \eqref{eq:two-field-curvature-measure} is
with respect to the support direction of a reconstructed spatial
set.  It introduces no new command direction and no derivative of
the observed response.  The same signed measure also gives the
entire surface area measure as its positive part.  When $C$ has
$C^2$ boundary and positive curvature, this part is
$r_C(\phi)\dd\phi$.

\begin{theorem}[Joint rigidity from two fixed forward fields]
\label{thm:two-field-rigidity}
Under \eqref{eq:two-field-gap}, the two fields
\eqref{eq:two-field-data} determine the complete canonical triple
\begin{equation}\label{eq:two-field-canonical-triple}
  \mathcal O_0=\mathcal O-s(A),\qquad
  A_0=A-s(A),\qquad \mu_0=(z\mapsto z-s(A))_\#\mu,
\end{equation}
with probability measures identified exactly.  In particular,
if $\mu$ has an $L^1$ density $j$, its centered density
$j_0(z)=j(z+s(A))$ is identified almost everywhere.
No angular continuum, short-flight limit, second stationary law,
or separately prescribed reciprocal preparation is required.
No periodicity assumption or lower bound on obstacle widths is
required.
\end{theorem}
\begin{proof}
Recover $v$ by Lemma~\ref{lem:two-field-prefix}, then take the
connected components $P$ of $\operatorname{supp}(v(x)\dd x)$.
Match the two collision
support components by Lemma~\ref{lem:two-field-supports}.
For each $P=P_C$, Lemma~\ref{lem:two-field-calibration} determines
$C$ up to translation, hence determines its Steiner-centered
representative $\overline C=C-s(C)$.
Set
\begin{equation}\label{eq:two-field-geometry-recovery}
                  C_P^0=\overline C+s(P).
\end{equation}
Steiner points add under Minkowski addition, so
$s(P)=s(C)-s(A)$.  Consequently $C_P^0=C-s(A)$, and
\[
 h_{-A_0}=h_P-h_{C_P^0},\qquad
                         \mathcal O_0=\bigcup_P C_P^0.
\]
Any one component gives the same centered footprint.

It remains to recover the probability law, which is not obtained merely
from its support.  Choose any $P=P_C$ and let $v_P=\one_Pv$,
extended by zero to the plane.  The positive component separation
shows that this is the occupation contributed by $C$ alone.
It is bounded, compactly supported, and integrable, and
\begin{equation}\label{eq:two-field-density-convolution}
                        v_P=\one_{C_P^0}*\check\mu_0,
 \qquad \check\mu_0=(z\mapsto-z)_\#\mu_0.
\end{equation}
With Fourier transform $\widehat f(\xi)=\int
e^{-i\xi\cdot x}f(x)\dd x$, one obtains
\begin{equation}\label{eq:two-field-density-quotient}
 \widehat\mu_0(-\xi)
             =\frac{\widehat v_P(\xi)}
                    {\widehat{\one_{C_P^0}}(\xi)}
 \quad\text{whenever }
                    \widehat{\one_{C_P^0}}(\xi)\ne0.
\end{equation}
The denominator is the restriction of an entire function, is
nonzero at the origin, and has a zero set with empty interior in
$\R^2$.  Indeed, vanishing on a real open rectangle would force
vanishing identically by applying the one-variable identity
theorem successively to the two coordinates.  The quotient
therefore specifies the continuous Fourier transform of $\mu_0$
on a dense set and hence everywhere.  Uniqueness of the Fourier
transform of a finite Borel measure recovers $\mu_0$ exactly.
All steps use only the two fields and the stated class bounds.
\end{proof}

\begin{corollary}[Complete fiber, periods, and response completion]
\label{cor:two-field-fiber}
Two admissible configurations and stationary laws have the same
ordered pair $(F_+,F_-)$ almost everywhere if and only if, for some
$h\in\R^2$,
\begin{equation}\label{eq:two-field-fiber}
 \mathcal O'=\mathcal O+h,\qquad A'=A+h,
                         \qquad \mu'=(z\mapsto z+h)_\#\mu.
\end{equation}
These conditions are equivalent to pointwise equality of every
physical forward mean at every nominal center and every finite
command displacement, with the prescribed closed-set bit convention.
Moreover,
\begin{equation}\label{eq:two-field-periods}
                  \operatorname{Per}(F_+,F_-)
                              =\operatorname{Per}(\mathcal O).
\end{equation}
This group is discrete, and it has rank two exactly when
$\mathcal O$ admits a rank-two lattice of periods with finitely
many obstacle orbits.  In \eqref{eq:two-field-periods}, field
periods may be defined almost everywhere or pointwise for the
physical representatives; both definitions give the same group.
Thus the two fixed fields determine both the full forward response
and the full obstacle period group.
\end{corollary}
\begin{proof}
Equality of the two fields gives equality of the canonical
triples by Theorem~\ref{thm:two-field-rigidity} and therefore
\eqref{eq:two-field-fiber}, with $h=s(A')-s(A)$.
Conversely, a common translation couples the actual starts,
segments, and obstacles and preserves every collision bit,
including its solid-start convention.  Recovery of the canonical
triple supplies the defining integral for every other command,
which proves the fiber and completion assertions.

If $w$ is an almost-everywhere common period of the two fields,
the finite prefix inverse shows that $v(\cdot+w)=v$ almost
everywhere.  Translation by $w$ consequently permutes the connected
components $P_C=C+(-A)$ of its measure support.
From $P_C+w=P_{C'}$, cancellation of their support functions gives
$C+w=C'$.  Hence $w$ is a period of $\mathcal O$.
The converse follows directly from translation covariance of the
bit.  Any nonzero period takes a compact obstacle to a different
one, so its length is at least $d$ and the group is discrete.
A discrete subgroup of $\R^2$ has rank at most two; if it has rank
two, local finiteness and bounded obstacle diameters imply that
only finitely many obstacle orbits meet a compact fundamental
cell.  The converse implication is immediate.
\end{proof}

The theorem concerns two exact spatial fields on the plane, and
their field values are still continuum data.  Its data reduction
is nevertheless literal: the command menu consists of the two
vectors $a,-a$, both of a fixed positive length.  The finite
prefix inverse already has a uniform finite-site stability bound.
Quantitative recovery of supports, the negative curvature atoms,
and the Fourier quotient from sampled bits requires the separate
finite estimates below.

\begin{proposition}[A single forward orientation]
\label{prop:two-field-one-orientation}
For any fixed compactly supported launch probability $\mu$, there
are distinct smooth strictly convex single-obstacle configurations
with the same forward response fields in direction $e_1$ at every
positive command length.  Their incoming arcs coincide and their
outgoing arcs differ.
\end{proposition}
\begin{proof}
Choose a nonzero smooth periodic function $b$ whose support is
contained in $\{\phi:\cos\phi>0\}$, and put
\[
 h_\lambda(\phi)=1+\lambda b(\phi),\qquad
               |\lambda|\|b+b''\|_\infty<1/2.
\]
Then $h_\lambda+h_\lambda''>1/2$, so $h_\lambda$ is the support
function of a smooth strictly convex body $C_\lambda$.
On the closed incoming normal semicircle $\cos\phi\le0$, both
$b$ and $b'$ vanish.  The support-point formula
$c_\lambda=h_\lambda n_\phi+h_\lambda' n_\phi^\perp$
therefore gives the same incoming semicircle for every $\lambda$.
Its transverse endpoints are $(0,\pm1)$, and the left endpoint
of every horizontal section at height $w\in[-1,1]$ is
$\ell(w)=-\sqrt{1-w^2}$, including the degenerate endpoint sections.

For any $T>0$, the closed collision convention gives exactly
\[
 B_{Te_1}^{C_\lambda}(u,w)
   =\one_{\{-1\le w\le1,\ \ell(w)-T\le u<\ell(w)\}}.
\]
A free start to the left hits if its segment reaches the left
endpoint; a solid start is zero; and a start to the right of the
section moves away.  The same description includes the two
tangent sections and their solid endpoints.  Thus the bits agree
pointwise for all $T>0$, and integration against the same $\mu$
gives identical response fields at every nominal center.
For $\lambda\ne0$, the support function differs from $h_0$ on
the outgoing normal semicircle.  The bodies are not translates:
a linear translation term vanishing on an open semicircle is zero,
whereas $\lambda b$ is not.  Their outgoing arcs therefore differ.
\end{proof}
